%%% intro.tex — Введение
\section*{ВВЕДЕНИЕ}
\addcontentsline{toc}{section}{ВВЕДЕНИЕ}

Современный этап развития информационных технологий характеризуется
стремительным ростом объёмов данных, передаваемых по компьютерным сетям,
и одновременным усложнением ландшафта киберугроз. По данным аналитических
отчётов ведущих компаний в области информационной безопасности, количество
инцидентов, связанных с утечками конфиденциальной информации, увеличивается
ежегодно на 15--25\,\%. При этом наблюдается устойчивая тенденция
к переходу злоумышленников от грубых методов атак к изощрённым техникам,
основанным на принципах скрытности и маскировки вредоносной активности.

Одним из наиболее сложных для обнаружения классов угроз является
использование стеганографических методов для скрытой передачи информации
по сетевым каналам. В отличие от криптографии, которая делает
содержимое сообщения нечитаемым, но не скрывает сам факт передачи,
стеганография направлена на сокрытие самого факта существования
коммуникации. Данное свойство делает стеганографию
привлекательным инструментом для организации скрытых каналов связи
в корпоративных сетях, реализации каналов управления вредоносным
программным обеспечением (C2-каналов) и эксфильтрации конфиденциальных
данных в обход систем глубокого анализа трафика (DPI), межсетевых
экранов и систем обнаружения вторжений (IDS).

Стек протоколов TCP/IP, лежащий в основе функционирования глобальной
сети Интернет и подавляющего большинства корпоративных сетей,
содержит значительное количество полей в заголовках
пакетов, которые могут быть использованы в качестве контейнеров для
встраивания скрытой информации. В частности, поле идентификации
IP-датаграммы (IP~Identification, 16~бит), начальный порядковый
номер TCP-сегмента (ISN, 32~бита), опция Timestamp протокола
TCP (32~бита поля TSval), область данных эхо-запроса ICMP
(до 1472~байт) и поле имени запроса DNS (QNAME, до 255~байт)
представляют собой потенциальные носители стеганографических вложений.
Каждое из указанных полей обладает различной ёмкостью, уровнем
скрытности и устойчивостью к обнаружению, что обуславливает
необходимость комплексного исследования методов их использования.

Существующие программные средства обнаружения вторжений, такие как
Zeek, Snort и Suricata, преимущественно
ориентированы на выявление известных сигнатур атак и не обладают
встроенными механизмами обнаружения стеганографических скрытых
каналов. Разработка специализированных детекторов, способных
выявлять аномалии в статистических характеристиках полей заголовков,
является актуальной научно-технической задачей.

\textbf{Объектом исследования} являются сетевые протоколы стека
TCP/IP как среда передачи данных в компьютерных сетях.

\textbf{Предметом исследования} выступают методы стеганографического
сокрытия данных в заголовках и полях сетевых пакетов, а также методы
обнаружения скрытых каналов передачи информации.

\textbf{Целью работы} является разработка и программная реализация
системы стеганографического сокрытия данных в сетевых пакетах
с криптографической защитой и многоуровневым модулем обнаружения
вредоносных вложений.

Для достижения поставленной цели сформулированы следующие
\textbf{задачи}:
\begin{enumerate}
  \item провести анализ современных методов сетевой стеганографии
        и подходов к обнаружению скрытых каналов;
  \item разработать архитектуру системы, включающей пять каналов
        сокрытия данных: ICMP payload, IP~ID, TCP~ISN, DNS~QNAME
        и TCP~Timestamp;
  \item реализовать криптографическую защиту передаваемых данных
        на основе алгоритма AES-256-GCM с функцией деривации ключа
        Argon2id и механизм фрагментации данных с контролем
        целостности;
  \item разработать модуль обнаружения скрытых каналов, объединяющий
        статистический, сигнатурный анализ и методы машинного
        обучения, а также модуль анализа содержимого для обнаружения
        вредоносных вложений;
  \item провести тестирование разработанной системы и выполнить
        анализ эффективности каждого канала сокрытия и метода
        обнаружения.
\end{enumerate}

\textbf{Методы исследования.} В работе применяются методы системного
анализа, теории информации (энтропия Шеннона, критерий хи-квадрат),
математической статистики (критерий Колмогорова--Смирнова),
машинного обучения (алгоритм Isolation Forest для обнаружения аномалий
и алгоритм Random Forest для supervised-классификации), сигнатурный
анализ содержимого payload, а также методы
объектно-ориентированного проектирования и программной инженерии.

\textbf{Практическая значимость} результатов работы заключается
в создании программного комплекса \texttt{NetStego}, который может
быть использован в следующих целях:
\begin{enumerate}[label=\arabic*)]
  \item исследование уязвимостей сетевых протоколов к
        стеганографическим атакам;
  \item обучение специалистов по информационной безопасности
        методам выявления скрытых каналов;
  \item тестирование и совершенствование корпоративных систем
        защиты от утечек данных.
\end{enumerate}

\textbf{Структура работы.} Выпускная квалификационная работа
состоит из введения, трёх глав, заключения, списка использованных
источников и пяти приложений.

В \textit{первой главе} рассматривается актуальность проблемы
стеганографической передачи данных в контексте современной
кибербезопасности: анализируется ландшафт киберугроз, описываются
примеры применения стеганографии злоумышленниками, исследуется
влияние скрытых каналов на безопасность корпоративных сетей
и обосновывается пригодность протоколов TCP/IP в качестве
носителей скрытой информации.

Во \textit{второй главе} изложены теоретические основы сетевой
стеганографии: классификация методов сокрытия, детальное описание
алгоритмов встраивания данных в заголовки пакетов, обзор методов
обнаружения скрытых каналов (статистических, сигнатурных и на
основе машинного обучения), а также сравнительный анализ
существующих программных решений.

В \textit{третьей главе} представлена программная реализация
системы \texttt{NetStego}: описаны используемые инструменты
разработки, архитектура системы, реализация каждого модуля
(криптографического ядра, каналов сокрытия, модуля обнаружения
скрытых каналов и анализа вредоносных вложений,
сетевых модулей отправки и приёма, системы сбора статистики
и визуализации), приведены результаты
тестирования производительности и эффективности обнаружения.

Общий объём работы составляет \pageref{LastPage}~страниц.
Список использованных источников включает 23~наименования.
Работа содержит 8~таблиц, 13~рисунков и 15~листингов.
